% Part: computability
% Chapter: computability-theory
% Section: halting-problem

\documentclass[../../../include/open-logic-section]{subfiles}

\begin{document}

% SEGMENT OLP-0235-S01

\olfileid{cmp}{thy}{hlt}
\olsection{நிற்றல் சிக்கல்}

அமைப்பின்படி, அனைத்தளாவிய பகுதியாகக் கணிக்கத்தக்க
சார்பு~$\fn{Un}(e,x)$ ஆனது, குறியீடு~$e$ கொண்ட சார்பின் கணக்கீடு
உள்ளீடு~$x$ க்கு ஒரு மதிப்பை உண்டாக்கினால், எனில் மட்டுமே,
வரையறுக்கப்பட்டுள்ளது. இது நிகழ்கிறதா என்று தீர்மானிக்க முடியுமா
எனக் கேட்பது இயல்பானது. உண்மையில் முடியாது. கணக்கீட்டின் டியூரிங்
பொறி மாதிரியைப் பொறுத்தவரை, கொடுக்கப்பட்ட டியூரிங் பொறி கொடுக்கப்பட்ட
உள்ளீட்டில் நிற்கிறதா என்பது கணிப்பியல் முறையில் தீர்மானிக்கவியலாதது.
எனவே பின்வரும் தேற்றம் “நிற்றல் சிக்கலின் தீர்மானிக்கவியலாமை”
என அறியப்படுகிறது. கீழே இரு நிரூபிப்புகளைத் தருவோம். முதலாவது நமது
முந்தைய உரையாடலின் இழையைத் தொடர்கிறது; இரண்டாவது மேலும் நேரடியானது.

\begin{thm}
\ollabel{thm:halting-problem}
பின்வருமாறு இருக்கட்டும்:
\[
h(e, x)  =
\begin{cases}
1 & \text{\/ $\fn{Un}(e, x)$ வரையறுக்கப்பட்டிருந்தால்} \\
0 & \text{இல்லையெனில்.}
\end{cases}
\]
அப்போது $h$ கணிக்கத்தக்கதல்ல.
\end{thm}

\begin{proof}
$h$ கணிக்கத்தக்கது என்று கொள்க. இது அனைத்தளாவிய கணிக்கத்தக்க சார்பு
ஒன்றை வரையறுக்க அனுமதிக்கும் என்று காட்டுவோம். பின்வருமாறு வரையறுக்க:
\[
\fn{Un'}(e,x) =
\begin{cases}
\fn{Un}(e,x) & \text{$h(e,x) = 1$ எனில்} \\
0 & \text{இல்லையெனில்.}
\end{cases}
\]
ஆனால் இப்போது $\fn{Un'}(e, x)$ ஒரு முழுச் சார்பு;~$h$
கணிக்கத்தக்கது எனில் அதுவும் கணிக்கத்தக்கது. எடுத்துக்காட்டாக,
முதன்மை மீள்வரையறையைப் பயன்படுத்தி $g$ ஐப் பின்வருமாறு
வரையறுக்கலாம்:
\begin{align*}
g(0, e, x) & \simeq 0 \\
g(y+1, e, x) & \simeq \fn{Un}(e,x);
\end{align*}
பின்னர்
\[
\fn{Un'}(e,x) \simeq g(h(e,x),e,x).
\]
பிந்தைய சார்பு வரையறுக்கப்பட்டுள்ள ஒவ்வோர் இடத்திலும்
$\fn{Un'}(e,x)$ ஆனது $\fn{Un}(e,x)$ உடன் உடன்படுவதால்,
முழுச் சார்புகளாக அமையும் பகுதியாகக் கணிக்கத்தக்க சார்புகளுக்காக
$\fn{Un'}$ அனைத்தளாவியது. ஆனால் இது
\olref[nou]{thm:no-univ} உடன் முரண்படுகிறது.
\end{proof}

\begin{proof}
$h(e,x)$ கணிக்கத்தக்கது என்று கொள்க. சார்பு $g$ ஐப் பின்வருமாறு
வரையறுக்க:
\[
g(x) =
\begin{cases}
  0                & \text{$h(x,x) = 0$ எனில்} \\
  \fundefined & \text{இல்லையெனில்.}
\end{cases}
\]
சார்பு $g$ பகுதியாகக் கணிக்கத்தக்கது. எடுத்துக்காட்டாக, அதை
$\umin{y}{h(x,x) = 0}$ என வரையறுக்கலாம். எனவே ஏதேனும்~$e$ க்கு,
ஒவ்வொரு~$x$ க்கும் $g(x) \simeq \fn{Un}(e, x)$. $g$ ஆனது $e$
இல் வரையறுக்கப்பட்டுள்ளதா? அவ்வாறெனில்,~$g$ இன் வரையறைப்படி
$h(e,e) = 0$ (~$h$ வரையறுக்கப்பட்டிருந்தால் மட்டுமே அது மதிப்பு~$0$
ஐ எடுக்க முடியும்).~$h$ இன் வரையறைப்படி, இது
$\fn{Un}(e, e)$ வரையறுக்கப்படவில்லை என்று பொருள்.
ஒவ்வொரு~$x$ க்கும் $g(x) \simeq \fn{Un}(e, x)$ என்ற நமது
அனுமானத்தின்படி, $g(e)$ வரையறுக்கப்படவில்லை; இது ஒரு முரண்பாடு.
மறுபுறம், $g(e)$ வரையறுக்கப்படவில்லை எனில் $h(e,e) \neq 0$;
எனவே $h(e,e) = 1$. இதிலிருந்து $\fn{Un}(e, e)$ வரையறுக்கப்பட்டுள்ளது
என்று கிடைக்கிறது. ஆனால் $g(x) \simeq \fn{Un}(e, x)$ என்பதால்,
$g(e)$ உம் வரையறுக்கப்பட்டிருக்கும். மீண்டும் ஒரு முரண்பாடு.
\end{proof}

\begin{tagblock}{TMs}
\begin{explain}
இவ்வாதத்தை டியூரிங் பொறிகளைக் கொண்டு விவரிக்கலாம். டியூரிங்
பொறி~$E$ ஒன்றின் விளக்கம் மற்றும் உள்ளீடு~$x$ ஒன்றை உள்ளீடாகப்
பெற்று, உள்ளீடு~$x$ இல் $E$ நிற்கிறதா இல்லையா என்று தீர்மானிக்கும்
டியூரிங் பொறி~$H$ ஒன்று உள்ளது என்று கொள்க. அப்போது தனி
உள்ளீடு~$x$ ஒன்றைப் பெற்று, சுட்டெண்~$x$ உடைய பொறி $M_x$ ஆனது
உள்ளீடு~$x$ இல் நிற்கிறதா என்று தீர்மானிக்க $H$ ஐ இயக்கி,
அதற்கு எதிர்மாறாகச் செய்யும் மற்றொரு டியூரிங் பொறி~$G$ ஐ அமைக்கலாம்.
வேறுவிதமாகக் கூறினால், உள்ளீடு~$x$ இல் $M_x$ நிற்கிறது என்று
$H$ அறிவித்தால் $G$ ஒரு முடிவிலிச் சுழற்சிக்குள் செல்கிறது;
உள்ளீடு~$x$ இல் $M_x$ நிற்பதில்லை என்று $H$ அறிவித்தால்,
$G$ வெறுமனே நிற்கிறது. தன் சுட்டெண்ணையே உள்ளீடாகக் கொண்டால்
$G$ நிற்கிறதா? அது நிற்காவிட்டால், எனில் மட்டுமே, நிற்கிறது என்று
மேலுள்ள வாதம் காட்டுகிறது---இது ஒரு முரண்பாடு. எனவே அத்தகைய
டியூரிங் பொறி~$H$ உள்ளது என்ற நமது அனுமானம் பொய்யாக வேண்டும்.
\end{explain}
\end{tagblock}

\end{document}
